\documentclass[12pt,a4paper]{article}

\usepackage[utf8]{inputenc}
\usepackage[T1]{fontenc}
\usepackage[french]{babel}
\usepackage{lmodern}
\usepackage{amsmath,amssymb}
\usepackage{geometry}
\usepackage{microtype}
\usepackage{hyperref}
\usepackage{enumitem}
\geometry{top=2.5cm,bottom=2.5cm,left=2.5cm,right=2.5cm}
\hypersetup{hidelinks,pdftitle={Documentation Consolidation HCANN}}

\title{Documentation du Module de Consolidation}
\author{Projet HCANN-Agent}
\date{\today}

\begin{document}
\maketitle

\section{Module Mémoire (\texttt{consolidation.py})}

\subsection{Objectif du module}
Ce module implémente la mémoire épisodique associative de HCANN au moyen d'un graphe hebbien dynamique:
\begin{itemize}
    \item \textbf{stockage d'épisodes} sous forme de n\oe uds indexés;
    \item \textbf{mise à jour des liens} par apprentissage hebbien avec décroissance;
    \item \textbf{récupération} via similarité cosinus + propagation d'activation;
    \item \textbf{consolidation} des hubs et oubli adaptatif.
\end{itemize}

\subsection{Composants principaux}
\begin{itemize}
    \item \texttt{HebbianMemoryGraph}: structure de données et opérations graphe;
    \item \texttt{HebbianConsolidation}: wrapper de haut niveau pour l'agent et l'entraînement.
\end{itemize}

\subsection{Vue théorique (graphe pondéré dynamique)}
Le système mémoire est modélisé par un graphe non orienté pondéré:
\begin{equation}
G_t = (V_t, E_t, W_t),
\end{equation}
où:
\begin{itemize}
    \item $V_t$ est l'ensemble des épisodes actifs à l'instant $t$;
    \item $E_t$ est l'ensemble des arêtes d'association;
    \item $W_t \in \mathbb{R}^{N\times N}$ contient les poids synaptiques (\texttt{edges}).
\end{itemize}

Chaque n\oe ud $i$ est associé à un état latent hippocampique $\mathbf{e}_i \in \mathbb{R}^{d_h}$.  
Le degré pondéré (force) d'un n\oe ud est:
\begin{equation}
\deg_w(i)=\sum_{j=1}^{N} W_{ij}.
\end{equation}
La détection des hubs est une approximation par degré binaire:
\begin{equation}
\deg_\theta(i)=\sum_{j=1}^{N}\mathbb{1}[W_{ij}>\theta], \quad
i \in \mathcal{H} \Leftrightarrow \deg_\theta(i) > \tau_h.
\end{equation}

\subsection{1) Structure du graphe}
Pour $N$ n\oe uds actifs:
\begin{itemize}
    \item matrice d'embeddings $E \in \mathbb{R}^{N \times d_h}$ (\texttt{node\_embeddings});
    \item matrice des poids hebbiens $W \in \mathbb{R}^{N \times N}$ (\texttt{edges});
    \item méta-données par n\oe ud: timestamp, nombre d'accès, drapeau \texttt{consolidated}.
\end{itemize}

\subsection{2) Mise à jour hebbienne}
La mise à jour des connexions suit:
\begin{equation}
W_{t+1} = \lambda W_t + \Delta W,
\end{equation}
où $\lambda=\texttt{hebbian\_decay}$ et $\Delta W$ ajoute un renforcement symétrique pour les paires co-activées:
\begin{equation}
\Delta W_{ij} = \Delta W_{ji} = \eta.
\end{equation}
Interprétation théorie$\rightarrow$calcul:
\begin{itemize}
    \item \textbf{théorie}: décroissance homogène de toutes les synapses (\textit{forgetting});
    \item \textbf{computation}: opération dense \texttt{self.edges *= hebbian\_decay};
    \item \textbf{théorie}: potentiel d'association entre épisodes co-activés;
    \item \textbf{computation}: \texttt{edges[idx1, idx2] += weight} et symétrique.
\end{itemize}

\subsection{3) Retrieval par spreading activation}
À partir d'une requête $\mathbf{q}$:
\begin{enumerate}[leftmargin=*]
    \item score de base par similarité cosinus:
    \begin{equation}
    s_i = \cos(\mathbf{q}, \mathbf{e}_i);
    \end{equation}
    \item propagation locale pour les n\oe uds activés:
    \begin{equation}
    \tilde{s}_j = s_j + \sum_{i \in \mathcal{A}} \alpha W_{ij},
    \end{equation}
    avec $\alpha=\texttt{spreading\_strength}$;
    \item sélection des \texttt{top-k} selon $\tilde{s}$.
\end{enumerate}

Sous forme matricielle, en notant $M$ un masque des n\oe uds actifs (score de base $>\theta_b$):
\begin{equation}
\tilde{\mathbf{s}} = \mathbf{s} + \alpha\,(M \odot \mathbf{s})\,W.
\end{equation}
Le code actuel implémente cette propagation de façon explicite par boucles locales sur les voisins non nuls.

\subsection{4) Détection de hubs et consolidation}
Un hub est un n\oe ud de degré supérieur à \texttt{hub\_threshold}.  
Le wrapper \texttt{HebbianConsolidation.consolidate}:
\begin{itemize}
    \item détecte les hubs;
    \item récupère leurs épisodes voisins;
    \item applique éventuellement un callback LLM pour produire un résumé;
    \item marque les hubs consolidés;
    \item déclenche l'oubli adaptatif.
\end{itemize}

Modélisation de la consolidation d'un hub $h$:
\begin{equation}
\mathcal{N}(h)=\{j \mid W_{hj}>\theta_e\}, \qquad
\text{summary}_h = f_{\text{LLM}}(\{x_j\}_{j \in \mathcal{N}(h)}),
\end{equation}
où $x_j$ représente les données épisodiques associées au n\oe ud $j$.

\subsection{5) Oubli adaptatif}
Un n\oe ud est candidat à suppression si:
\begin{equation}
(\text{poids total} < \tau_w)\ \land\ (\text{âge} > \tau_t)\ \land\ (\text{accès} = 0),
\end{equation}
avec $\tau_w=\texttt{prune\_weight\_thresh}$ et $\tau_t=\texttt{prune\_age\_thresh}$.

\subsection{Rendu computationnel (algorithmes)}
\paragraph{Algorithme A: insertion d'un épisode}
\begin{enumerate}[leftmargin=*]
    \item Vérifier capacité $N<C$ sinon lancer l'oubli adaptatif;
    \item écrire l'embedding dans \texttt{node\_embeddings[idx]};
    \item initialiser timestamp et compteur d'accès;
    \item mettre à jour dictionnaires d'indexation.
\end{enumerate}
Coût: $O(1)$ amorti hors oubli.

\paragraph{Algorithme B: retrieval 2 phases}
\begin{enumerate}[leftmargin=*]
    \item Calcul cosinus requête$\leftrightarrow$tous les n\oe uds actifs;
    \item pour chaque n\oe ud activé, parcourir ses voisins et booster les scores;
    \item extraire \texttt{top-k} via \texttt{torch.topk}.
\end{enumerate}
Coût dominant:
\begin{equation}
O(Nd_h) + O\!\left(\sum_{i \in \mathcal{A}} |\mathcal{N}(i)|\right) + O(N \log k).
\end{equation}

\paragraph{Algorithme C: consolidation périodique}
\begin{enumerate}[leftmargin=*]
    \item Détecter les hubs par degrés seuilés;
    \item agréger les voisins des hubs non consolidés;
    \item produire un résumé (optionnel) puis marquer consolidé;
    \item appliquer l'oubli adaptatif.
\end{enumerate}

\subsection{Correspondance théorie $\leftrightarrow$ code}
\begin{itemize}
    \item \textbf{objet théorique} $G_t$ $\rightarrow$ \texttt{nodes + edges + id\_to\_index};
    \item \textbf{poids synaptiques} $W$ $\rightarrow$ \texttt{self.edges};
    \item \textbf{degré / hubs} $\deg_\theta$ $\rightarrow$ \texttt{detect\_hubs()};
    \item \textbf{activation propagée} $\tilde{\mathbf{s}}$ $\rightarrow$ \texttt{spreading\_activation()};
    \item \textbf{stabilité/plasticité} $(\lambda,\eta)$ $\rightarrow$ decay global + incréments locaux.
\end{itemize}

\subsection{Correspondance implémentation}
\begin{itemize}
    \item \texttt{add\_node}: insertion d'un épisode et indexation;
    \item \texttt{update\_hebbian\_weights}: decay global puis renforcement pairwise;
    \item \texttt{spreading\_activation}: scoring 2 phases (direct + propagation);
    \item \texttt{detect\_hubs}: calcul de degré sur seuil de connexion;
    \item \texttt{adaptive\_forgetting}: suppression multi-critères;
    \item \texttt{consolidate}: orchestration consolidation + oubli.
\end{itemize}

\subsection{Interface détaillée}

\paragraph{\texttt{HebbianMemoryGraph.add\_node(node\_id, data, embedding)}}
\begin{itemize}[leftmargin=*]
    \item \textbf{Entrées}: identifiant n\oe ud, dictionnaire de données, embedding NumPy $(d_h)$.
    \item \textbf{Effet}: ajoute le n\oe ud (ou déclenche un oubli adaptatif si capacité atteinte).
\end{itemize}

\paragraph{\texttt{HebbianMemoryGraph.update\_hebbian\_weights(node\_id\_1, node\_id\_2, weight)}}
\begin{itemize}[leftmargin=*]
    \item \textbf{Entrées}: deux identifiants et une force de renforcement.
    \item \textbf{Effet}: applique décroissance globale puis renforcement symétrique des arêtes.
\end{itemize}

\paragraph{\texttt{HebbianMemoryGraph.spreading\_activation(query\_embedding, k=5)}}
\begin{itemize}[leftmargin=*]
    \item \textbf{Entrées}: requête NumPy $(d_h)$, nombre de résultats.
    \item \textbf{Sortie}: liste d'épisodes enrichis (\texttt{id}, \texttt{score}, \texttt{data}, etc.).
\end{itemize}

\paragraph{\texttt{HebbianConsolidation.retrieve(query\_embedding, k=5)}}
\begin{itemize}[leftmargin=*]
    \item \textbf{Rôle}: alias haut niveau vers \texttt{spreading\_activation}.
\end{itemize}

\paragraph{\texttt{HebbianConsolidation.consolidate(llm\_callback=None)}}
\begin{itemize}[leftmargin=*]
    \item \textbf{Entrée}: callback optionnel prenant une liste d'épisodes.
    \item \textbf{Sortie}: liste des identifiants de hubs consolidés.
\end{itemize}

\paragraph{\texttt{HebbianConsolidation.get\_stats()}}
\begin{itemize}[leftmargin=*]
    \item \textbf{Sortie}: dictionnaire de métriques (\texttt{total\_nodes}, \texttt{total\_edges}, \texttt{hubs}, \texttt{consolidated}).
\end{itemize}

\subsection{Résumé opérationnel}
\begin{enumerate}[leftmargin=*]
    \item Ajouter un épisode avec son embedding hippocampique;
    \item renforcer les connexions avec les épisodes liés;
    \item récupérer les souvenirs pertinents par propagation;
    \item consolider périodiquement les hubs;
    \item élaguer les traces obsolètes pour contenir la mémoire.
\end{enumerate}

\end{document}
